\subsection{The Temperature ($\tau$) Effect}
\label{sec:tau-discussion}

Our central finding on the tabular side is that the prediction
soft-max temperature $\tau$ is the dominant knob for whether DRO
wins or loses on Adult DP.  Under a stepped schedule
($\tau{=}100$ for $\alpha\le 0.3$, $\tau{=}1$ for $\alpha{=}0.4$)
DRO loses on Adult DP almost everywhere; under a fixed $\tau{=}1$
schedule DRO wins on Adult DP at $\alpha\le0.2$ (no method claim is
made at $\alpha\ge0.3$ on Adult and Credit, where both methods fall
below the constant predictor; LSAC/DP is pinned at majority accuracy).
At $\alpha\le0.2$ the gap widens monotonically.
This is consistent with the theory: $\tau{=}100$
sharpens the predictions toward $0$/$1$ and so makes the inner
re-weighting concentrate weights on a few ``worst'' samples,
driving $\lambda_{\mathrm{DP}}$ to its clamp and causing the
$\lambda$-runaway the earlier ``adversarial feedback loop''
discussion was observing.  At $\tau{=}1$ the soft predictions
distribute weight smoothly and $\lambda$ stays bounded, so the
re-weighting helps rather than hurts.

\subsection{Adversarial vs.\ Random Corruption}
A historical pilot (pre-canonical protocol) found adversarial
fairness-targeted PGD raises Adult/Credit DP far more than matched-$\alpha$
random label noise.  We keep that qualitative sanity check (adversarial
$\gg$ random) but do \emph{not} treat older multiplicative factors as locked
canonical numbers until a seed-paired re-run under the $\tau{=}1$,
$k_{\mathrm{inner}}{=}10$, $n{=}6$ protocol is committed.  This still
rules out a pure ``attack is too weak'' reading of the Adult low-$\alpha$
cells under PGD.

\subsection{Mechanism (old ``feedback loop'')}
The earlier ``adversarial feedback loop'' account of Adult
$\alpha\ge 0.3$ (DRO collapsing to constant predictions) is no
longer the right diagnosis under $\tau{=}1$.  We do still see lower
accuracy at $\alpha{=}0.4$ ($\approx 0.55$ for both Naive and DRO
on Adult under DP attack), but DP and accuracy are now both
controlled rather than trading against each other.

\subsection{Empirical Radii (Q5)}
The TV-radius formula
$\rho_{\mathrm{DP},j}=\alpha/((1-\alpha)\pi_j+\alpha)$ and the
bias-corrected $\hat\pi_{j,\mathrm{clean}}=(\hat\pi_j-\alpha)/(1-2\alpha)$
are empirical, not theoretical.  Under our coordinated (70\%-minority)
attack structure they can be tightened by estimating $\pi_{\mathrm{clean}}$
from the observed post-attack group proportions directly (no per-sample mask; only $\alpha$ + known 70/30 structure).  We
implement this as \texttt{radii\_mode='empirical'} in
\texttt{src/training/dro\_fair.py} (default stays 'uniform').  This
is in line with Kuldeep's Q5 ("if the attack is known then we can
use this approximation according to the attack") and does not claim
the paper is wrong.
See report appendix for the clean inversion derivation: $\pi_{\mathrm{clean}}[\mathrm{min}] = \pi_{\mathrm{obs}}[\mathrm{min}] + 0.4\alpha$, etc. (B's math).

\subsection{Image data (UTKFace) --- real pilot; not a tabular copy}
The earlier ``DRO inverts on ResNet18 features'' claim is withdrawn
(synthetic smoke only).  The full 90-config REAL grid
(\S\ref{sec:results-utkface}) shows a \textbf{mixed} clean-test pattern:
stronger DRO DP wins mainly at high $\alpha$, not the Adult/Credit
$\alpha\le0.2$ regime.  Image features and tabular data should not be
conflated in the main claim.

\subsection{Limitations}
\label{sec:limitations}
We report three honest limitations that bound the scope of the claims
above.
\begin{enumerate}
\item \textbf{LSAC/DP degeneracy.}  Under the DP attack on LSAC, the
DRO model collapses to the constant (majority-class) predictor:
accuracy is pinned to the LSAC majority rate $\approx0.9016$ at every
$\alpha$, and DP does not move monotonically with corruption (it peaks
at $\alpha{=}0.1$ and falls at $\alpha\ge0.2$).  DRO loses to Naive on
LSAC/DP at every $\alpha$ ($0/6$ seeds).  This is an artifact of the TV
radii formula $\rho_{\mathrm{DP},j}=\alpha/((1-\alpha)\pi_j+\alpha)$ on
LSAC's $\sim90/10$ imbalanced protected group (the minority radius
blows up and over-weights a near-empty group), \emph{not} evidence that
``DRO loses on LSAC.''  Full diagnosis in
\texttt{docs/LSAC\_DEGENERACY.md}; LSAC/DP is reported as a degenerate
diagnostic and excluded from the DP win table, while LSAC/Combined is a
genuine clean win ($p{=}0.0156$ at $\alpha\in\{0.1,0.3,0.4\}$).

\item \textbf{$\alpha\ge0.3$ constant-predictor bound (Adult+Credit).}
At $\alpha\ge0.3$ on Adult and Credit, both methods' accuracy falls
below the per-dataset constant-predictor baseline (Adult
$\approx0.7521$, Credit $\approx0.7788$).  We make no method claim in
this regime, even when fairness gaps still favour DRO.  LSAC is not
``below'' its baseline --- it is pinned \emph{at} it ($\approx0.9016$)
under the DP attack, which is the degeneracy of item~1.

\item \textbf{IF story is MIXED.}  The IF \emph{metric} supports DRO on
Adult/Credit at $\alpha\in\{0.1$--$0.4\}$ including $\alpha{=}0.3$
($6/6$, $p{=}0.0156$; \S\ref{sec:results-tabular}), but
DP-under-IF-attack is not a second clean DP story: Adult $\alpha{=}0.3$
loses on DP ($\mathrm{DP}_{\mathrm{DRO}}{=}0.0241 >
\mathrm{DP}_{\mathrm{Naive}}{=}0.0227$, $1/6$) even while IF wins, and
LSAC/IF loses on DP for $\alpha\le0.3$.  The IF and DP metrics are
\emph{coupled}, not interchangeable; we state the IF win clearly but
pair it with the DP-coupling caveat rather than claiming a three-attack
DP sweep.
\end{enumerate}
